% 第二部分：引言与相关工作（由 main_zh.tex 引入）。

\section{引言}\label{sec:intro}

基于可验证奖励的强化学习（RLVR）在答案可以检验的问题上训练语言模型，GRPO~\cite{shao2024deepseekmath,deepseek2025r1} 及其变体~\cite{liu2025understanding,liu2026tricks}是最常见的形式。验证器精确但粗糙：它只判断最终答案对不对，不管推理是否站得住；当一组样本全对或全错时，它不给任何信号。评判模型可以直接给推理打分，过程监督~\cite{lightman2024verify,wang2024mathshepherd}和生成式验证器~\cite{zhang2024genverifier}是这一想法的两种形式。评判模型会带来显存和时间上的代价，也带来风险：学到的奖励可能被策略利用~\cite{gao2023scaling,skalse2022defining}。

对于 Qwen3.5-4B 策略的 RLVR，\kev{}~\cite{kev4b} 是一个加入成本格外低的评判模型。它是同一基座模型 Qwen3.5-4B-Base 上的一个 LoRA 适配器加一个指针头，因此策略和评判模型可以共用一份骨干权重：策略 LoRA 与评判 LoRA 并列挂在骨干上，评判模型一次前向传播就读出答案，无需生成文本。本文测量这种设计节省了什么，以及评判模型给 GRPO 带来了什么。我们提出三个问题：
\begin{enumerate}
\item 在同一进程、同一训练步下测量，共享骨干比把评判模型作为第二个模型加载节省多少显存（\S\ref{sec:res-shared}）？
\item 与同家族、同大小的 AR 评判模型相比，以及与该模型先推理再作答相比，\kev{} 区分策略解答正误的能力如何？每个评判模型有多容易被不该影响它的改动所左右（\S\ref{sec:res-study}）？
\item 把 \kev{} 的判断加入验证器奖励，能否改进 MATH 上的 GRPO？\kev{} 能否取代验证器（\S\ref{sec:res-grpo}）？
\end{enumerate}

\section{相关工作}\label{sec:related}

\paragraph{RLVR 与 GRPO。} GRPO~\cite{shao2024deepseekmath} 用同一提示下多个样本的组基线取代 PPO 的价值函数，DeepSeek-R1~\cite{deepseek2025r1} 使其成为推理训练的标准做法。后续分析表明细节很重要：长度归一化和组内标准差~\cite{liu2025understanding}，以及优势尺度的选择~\cite{liu2026tricks}。我们采用组均值基线和 \citet{liu2026tricks} 推荐的批级尺度、词元级损失，不加 KL 项。现代 RL 框架用推理引擎采样、用另一套框架训练，两者对同一组权重算出的概率略有差异，截断重要性采样可以修正这一点~\cite{yao2025offpolicy}。

\paragraph{用于推理的评判模型和奖励模型。} 结果奖励模型和过程奖励模型分别给解答或步骤打分~\cite{lightman2024verify,wang2024mathshepherd}。生成式验证器~\cite{zhang2024genverifier}询问语言模型某个解答是否正确（可以先推理），并读取“是”的概率。已知 LLM 评判对长度、风格和“声称正确”的说法很敏感~\cite{zheng2023judging}。策略被推得越狠，对奖励模型的过度优化就越严重~\cite{gao2023scaling}；在许多场景中，RL 策略都学会了利用其奖励信号~\cite{skalse2022defining,denison2024sycophancy,baker2025monitoring}。

\paragraph{策略与奖励之间共享权重。} 多 LoRA 推理~\cite{hu2022lora,kwon2023vllm}让多个适配器在推理时共用一个基座模型。这里把同一机制用在训练器内部：一个正在训练的策略适配器和一个冻结的评判适配器。使之严格成立的条件是，训练从不改变评判模型训练时所用的基座权重。
